% !TEX program = lualatex
% Transparencias de Fundamentos de las Matemáticas I
% Clase 11: Relaciones de equivalencia y de orden

\documentclass[aspectratio=169,11pt]{beamer}

\usepackage{fontspec}
\usepackage[spanish,es-nodecimaldot]{babel}
\usepackage{amsmath,amssymb,mathtools}
\usepackage{booktabs}
\usepackage{csquotes}
\usepackage{xcolor}

\usetheme{Madrid}
\usecolortheme{default}
\setbeamertemplate{navigation symbols}{}
\setbeamertemplate{footline}[frame number]

\definecolor{FdMBlue}{RGB}{20,55,100}
\definecolor{FdMRed}{RGB}{130,30,30}
\definecolor{FdMGreen}{RGB}{25,100,60}

\setbeamercolor{structure}{fg=FdMBlue}
\setbeamercolor{title}{fg=white,bg=FdMBlue}
\setbeamercolor{frametitle}{fg=white,bg=FdMBlue}
\setbeamercolor{block title}{fg=white,bg=FdMBlue}
\setbeamercolor{block body}{bg=blue!3}
\setbeamercolor{alerted text}{fg=FdMRed}

\setmainfont{Latin Modern Roman}
\setsansfont{Latin Modern Sans}
\setmonofont{Latin Modern Mono}

\newcommand{\N}{\mathbb{N}}
\newcommand{\Z}{\mathbb{Z}}
\newcommand{\Q}{\mathbb{Q}}
\newcommand{\R}{\mathbb{R}}
\newcommand{\C}{\mathbb{C}}
\newcommand{\Pow}{\mathcal{P}}
\newcommand{\id}{\operatorname{id}}
\newcommand{\mcd}{\operatorname{mcd}}
\newcommand{\mcm}{\operatorname{mcm}}

\title[FdM I -- Clase 11]{Relaciones de equivalencia y de orden}
\subtitle{Clasificar y comparar}
\author{Antonio Falcó Montesinos}
\institute{Universidad CEU Cardenal Herrera}
\date{Fundamentos de las Matemáticas I}

\begin{document}

\begin{frame}
  \titlepage
\end{frame}

\begin{frame}{Objetivos de la clase}
\begin{itemize}
  \item Definir relaciones de equivalencia y clases de equivalencia.
  \item Definir relaciones de orden.
  \item Distinguir clasificar de comparar.
\end{itemize}
\end{frame}

\begin{frame}{Mapa de la clase}
\tableofcontents
\end{frame}

\section{Equivalencia}

\begin{frame}{Equivalencia}
\begin{block}{Relación de equivalencia}
\begin{itemize}
  \item Una relación es de equivalencia si es reflexiva, simétrica y transitiva.
  \item Sirve para agrupar objetos que se consideran equivalentes bajo un criterio.
  \item Ejemplo: tener el mismo resto al dividir por \(n\).
\end{itemize}
\end{block}
\end{frame}

\begin{frame}{Equivalencia}
\begin{block}{Clases de equivalencia}
\begin{itemize}
  \item La clase de \(a\) es \([a]=\{x\in A:x\sim a\}\).
  \item Las clases agrupan todos los elementos equivalentes a uno dado.
  \item Dos clases son iguales o disjuntas.
\end{itemize}
\end{block}
\end{frame}

\begin{frame}{Equivalencia}
\begin{block}{Partición}
\begin{itemize}
  \item Una relación de equivalencia divide el conjunto en bloques.
  \item Cada elemento pertenece a exactamente una clase.
  \item Esta idea será esencial para construir cocientes.
\end{itemize}
\end{block}
\end{frame}

\section{Orden}

\begin{frame}{Orden}
\begin{block}{Relación de orden}
\begin{itemize}
  \item Una relación de orden es reflexiva, antisimétrica y transitiva.
  \item Permite comparar elementos de un conjunto.
  \item Si todos los pares son comparables, el orden es total.
\end{itemize}
\end{block}
\end{frame}

\begin{frame}{Orden}
\begin{block}{Comparación con equivalencia}
\begin{itemize}
  \item Equivalencia: agrupa objetos.
  \item Orden: compara objetos.
  \item Ambas se definen mediante propiedades lógicas de una relación.
\end{itemize}
\end{block}
\end{frame}

\section{Profundización}

\begin{frame}{Profundización}
\begin{block}{Clases: idea operativa}
\begin{itemize}
  \item La clase \([a]\) contiene todos los elementos indistinguibles de \(a\) según la relación.
  \item Cambiar el representante no cambia la clase si los representantes son equivalentes.
  \item Esta idea anticipa fracciones, congruencias y construcciones cociente.
\end{itemize}
\end{block}
\end{frame}

\begin{frame}{Profundización}
\begin{block}{Orden parcial y total}
\begin{itemize}
  \item En un orden parcial puede haber elementos no comparables.
  \item En un orden total, cualesquiera dos elementos son comparables.
  \item \(\subseteq\) en \(\mathcal{P}(A)\) es un orden parcial.
\end{itemize}
\end{block}
\end{frame}

\begin{frame}{Profundización}
\begin{block}{Comparación conceptual}
\begin{itemize}
  \item Equivalencia: sustituir un objeto por otro de la misma clase.
  \item Orden: decidir precedencia, contención o tamaño relativo.
  \item Ambas nociones dependen de propiedades lógicas precisas.
\end{itemize}
\end{block}
\end{frame}

\begin{frame}{Cierre}
\begin{itemize}
  \item Las equivalencias producen clases.
  \item Los órdenes permiten comparar.
  \item Ambas nociones preparan estructuras más ricas.
\end{itemize}
\end{frame}

\begin{frame}{Trabajo recomendado}
\begin{itemize}
  \item Releer las secciones correspondientes del manual.
  \item Identificar definiciones, símbolos nuevos y enunciados principales.
  \item Preparar dudas y ejemplos para el seminario asociado.
\end{itemize}
\end{frame}

\end{document}
